% ============================================================
% Section 4 Supplement: Quantitative Error Bounds
% Addresses remaining gaps from Codex review:
%   1. R_n quantitative bound
%   2. A2 (conditional independence) quantitative error
%   3. Multi-layer noise accumulation
%   4. P_e consistency with I
% ============================================================

\subsection{Quantitative Error Bounds}
\label{subsec:quantitative_errors}

We now provide the quantitative error bounds that complete the rigorous framework.

\subsubsection{Explicit Assumption Regime}

\begin{assumption}[Validity Regime for SNR Framework]
\label{ass:validity_regime}
The theoretical results in this section hold under the following explicit conditions:
\begin{enumerate}
    \item[(R1)] \textbf{Sparse graph regime}: Average degree $k = O(1)$ as $n \to \infty$. More precisely, we require $k = o(n^{1/3})$ to ensure the local tree approximation holds.

    \item[(R2)] \textbf{Layer depth constraint}: The number of aggregation layers satisfies
    \begin{equation}
        t \leq \left(\frac{1}{2} - \epsilon\right) \log_{k-1} n
    \end{equation}
    for some fixed $\epsilon > 0$. This ensures the cycle contamination term $(k-1)^{2t}/n \to 0$.

    \item[(R3)] \textbf{Feature dimension}: The feature dimension $d_x$ is fixed (does not grow with $n$).

    \item[(R4)] \textbf{Moderate feature SNR}: The feature signal-to-noise ratio satisfies $\|\mu\|^2/\sigma^2 = O(1)$, i.e., features are informative but not perfectly separating. This also implies $\|\mu\| = O(1)$ is bounded uniformly in $n$.

    \item[(R5)] \textbf{Non-degenerate correlation}: The edge correlation $|\rho| = |2h-1|$ is bounded away from 0 and 1, i.e., $\rho \in [\rho_{\min}, \rho_{\max}]$ for fixed constants $0 < \rho_{\min} < \rho_{\max} < 1$.

    \item[(R6)] \textbf{Positive definite covariance}: The noise covariance $\Sigma \succ 0$ is strictly positive definite with bounded condition number $\kappa(\Sigma) := \|\Sigma\|_{\mathrm{op}} \|\Sigma^{-1}\|_{\mathrm{op}} = O(1)$. This ensures the discriminability $\mathcal{D}(\mu, \Sigma) = 4\mu^\top \Sigma^{-1} \mu$ is well-defined and the Lipschitz bounds are finite.

    \item[(R7)] \textbf{Homoscedastic noise}: The conditional covariance $\mathrm{Var}(X_v \mid Y_v) = \Sigma$ is identical for both classes $Y_v = \pm 1$. This ``equal within-class variance'' assumption is standard in linear discriminant analysis and ensures $\Delta\Sigma = 0$ in our perturbation analysis.
\end{enumerate}
\end{assumption}

\begin{remark}[Constant Dependencies]
\label{rem:constant_dependencies}
The universal constants $C, C_1, C_2, C_3$ appearing in our bounds depend on:
\begin{itemize}
    \item The SBM parameters $(a, b)$ through $\rho = (a-b)/(a+b)$
    \item The noise covariance $\|\Sigma\|_{\mathrm{op}}$ and $\|\Sigma^{-1}\|_{\mathrm{op}}$
    \item The feature dimension $d_x$
    \item The sub-Gaussian constant of the noise distribution (if applicable)
\end{itemize}
These constants are \textbf{independent of} $n$ (graph size) and $t$ (layer depth), which is crucial for the asymptotic validity of our results.

\textbf{Uniformity clarification:} When we write $\kappa(\Sigma) = O(1)$, $\|\Sigma^{-1}\|_{\mathrm{op}} = O(1)$, or $\|\mu\| = O(1)$, these are \textit{uniform bounds} as $n \to \infty$ with all other model parameters $(\mu, \Sigma, a, b, d_x)$ held \textit{fixed}. The feature dimension $d_x$ is treated as a constant (Assumption R3), so these bounds do not involve high-dimensional asymptotics. If one wishes to consider $d_x \to \infty$ simultaneously, additional assumptions on the scaling of $\|\mu\|$ and $\|\Sigma\|$ with $d_x$ would be required.
\end{remark}

\subsubsection{Remainder Term $R_n$ Bound}

\begin{lemma}[Spectral Norm Bound on $R_n$]
\label{lem:R_n_bound}
Under the sparse CSBM (Definition~\ref{def:sparse_csbm}), the remainder term $R_n$ in the variance decomposition (Theorem~\ref{thm:agg_discriminability_rigorous}(b)) satisfies:
\begin{equation}
    \|R_n\|_{\mathrm{op}} \leq C \cdot \frac{k^2}{n} \cdot \|\Sigma\|_{\mathrm{op}}
\end{equation}
for a universal constant $C > 0$, with probability at least $1 - O(n^{-2})$.

Consequently, the relative error in discriminability is:
\begin{equation}
    \left|\frac{\mathcal{D}(A_v)}{\rho^2 k_{\mathrm{eff}} \mathcal{D}(X)} - 1\right| \leq \frac{C k^2}{n} + \frac{\|\mu\|^2 (1-\rho^2)}{\sigma^2 k_{\mathrm{eff}}}
\end{equation}
\end{lemma}

\begin{proof}
The remainder $R_n$ arises from correlations between neighbor labels in finite graphs. In the sparse CSBM, for any two neighbors $u, u' \in \mathcal{N}(v)$:
\begin{equation}
    \mathrm{Cov}(Y_u, Y_{u'} \mid Y_v) = O(k/n)
\end{equation}
This is because $u$ and $u'$ share a common neighbor (namely $v$) with probability $O(k/n)$ in the sparse regime.

Summing over all $O(k^2)$ pairs of neighbors:
\begin{equation}
    \|R_n\|_{\mathrm{op}} \leq k^2 \cdot O(k/n) \cdot \|\mu \mu^\top\|_{\mathrm{op}} = O(k^3 \|\mu\|^2 / n)
\end{equation}

For the discriminability, the error propagates as:
\begin{align}
    \mathcal{D}(A_v) &= \Delta_A^\top \Sigma_A^{-1} \Delta_A \\
    &= \Delta_A^\top \left(\Sigma_A^{(0)} + R_n\right)^{-1} \Delta_A \\
    &= \Delta_A^\top (\Sigma_A^{(0)})^{-1} \Delta_A \cdot (1 + O(\|R_n\|_{\mathrm{op}} / \|\Sigma_A^{(0)}\|_{\mathrm{op}}))
\end{align}
where $\Sigma_A^{(0)} = (\sum_u w_{vu}^2) \Sigma$ is the leading term.
\end{proof}

\subsubsection{Quantitative Conditional Independence (A2)}

\begin{lemma}[Total Variation Bound for Neighbor Independence]
\label{lem:A2_quantitative}
In the sparse CSBM with average degree $k = O(1)$, let $\mathcal{N}(v) = \{u_1, \ldots, u_d\}$ be the neighbors of node $v$. Define:
\begin{itemize}
    \item $P_{\mathrm{joint}}$: The true joint distribution of $(Y_{u_1}, \ldots, Y_{u_d}) \mid Y_v$
    \item $P_{\mathrm{prod}}$: The product distribution $\prod_{i=1}^d P(Y_{u_i} \mid Y_v)$
\end{itemize}

Then:
\begin{equation}
    d_{\mathrm{TV}}(P_{\mathrm{joint}}, P_{\mathrm{prod}}) \leq \frac{C k^4}{n}
\end{equation}
for a universal constant $C \leq 8$. Under assumption (R1) with $k = O(1)$, this is $O(1/n)$.
\end{lemma}

\begin{proof}[Proof via Coupling]
We construct an explicit coupling between $P_{\mathrm{joint}}$ and $P_{\mathrm{prod}}$.

\textbf{Step 1 (Bad event definition):} Define the ``bad event'':
\begin{equation}
    \mathcal{B} := \bigcup_{i < j} \mathcal{B}_{ij}, \quad \text{where } \mathcal{B}_{ij} := \{u_i \text{ and } u_j \text{ share a neighbor other than } v\}
\end{equation}

\textbf{Step 2 (Bad event probability --- detailed derivation):} We compute $\mathbb{P}(\mathcal{B})$ rigorously.

\textit{Step 2a (Per-pair probability):} Fix a pair of neighbors $(u_i, u_j) \in \mathcal{N}(v)^2$ with $i \neq j$. The event $\mathcal{B}_{ij}$ requires the existence of some node $w \neq v$ adjacent to both $u_i$ and $u_j$. By union bound over candidate nodes $w$:
\begin{align}
    \mathbb{P}(\mathcal{B}_{ij}) &\leq \sum_{w \neq v, u_i, u_j} \mathbb{P}((u_i, w) \in E) \cdot \mathbb{P}((u_j, w) \in E) \\
    &\leq (n - 3) \cdot \frac{a+b}{n} \cdot \frac{a+b}{n} \\
    &= (n-3) \cdot \frac{(a+b)^2}{n^2} = \frac{(a+b)^2 (n-3)}{n^2}
\end{align}
Since $k = (a+b)/2$, we have $(a+b) = 2k$, so:
\begin{equation}
    \mathbb{P}(\mathcal{B}_{ij}) \leq \frac{4k^2 (n-3)}{n^2} \leq \frac{4k^2}{n} \quad \text{for } n \geq 4
\end{equation}

\textit{Step 2b (Union bound over pairs --- high probability version):} The number of neighbor pairs is $\binom{|\mathcal{N}(v)|}{2}$.

\textbf{Degree distribution justification:} In the sparse CSBM (Definition~\ref{def:sparse_csbm}), the degree $|\mathcal{N}(v)|$ follows $\mathrm{Bin}(n-1, (a+b)/n)$ marginally. In the sparse limit where $k = (a+b)/2 = O(1)$ is fixed as $n \to \infty$, by the Poisson limit theorem (Le Cam's inequality), this converges in distribution to $\mathrm{Pois}(k)$ with total variation error $O(k^2/n)$. Thus, Poisson tail bounds apply asymptotically. For finite $n$, one may alternatively use Binomial Chernoff bounds directly, which yield the same exponential decay.

For Poisson$(k)$ degree distribution, by standard tail bounds (Chernoff or sub-Poisson concentration):
\begin{equation}
    \mathbb{P}(|\mathcal{N}(v)| > 2k) \leq e^{-k/3} \quad \text{for } k \geq 1
\end{equation}
Conditioning on the event $\{|\mathcal{N}(v)| \leq 2k\}$, which holds with probability $\geq 1 - e^{-k/3}$:
\begin{equation}
    \mathbb{P}(\mathcal{B} \mid |\mathcal{N}(v)| \leq 2k) \leq \binom{2k}{2} \cdot \frac{4k^2}{n} = \frac{2k(2k-1)}{2} \cdot \frac{4k^2}{n} \leq \frac{8k^4}{n}
\end{equation}
By law of total probability:
\begin{equation}
    \mathbb{P}(\mathcal{B}) \leq \frac{8k^4}{n} + e^{-k/3} = O\left(\frac{k^4}{n}\right) + O(e^{-k/3})
\end{equation}
For $k = O(1)$ fixed and $n \to \infty$, the first term dominates, giving $\mathbb{P}(\mathcal{B}) = O(k^4/n)$.

\textit{Step 2c (Alternative: expectation-based bound):} Taking expectation over $|\mathcal{N}(v)|$ directly:
\begin{equation}
    \mathbb{E}\left[\binom{|\mathcal{N}(v)|}{2}\right] = \frac{\mathbb{E}[|\mathcal{N}(v)|^2] - \mathbb{E}[|\mathcal{N}(v)|]}{2} \leq \frac{k^2 + k - k}{2} = \frac{k^2}{2}
\end{equation}
(using $\mathbb{E}[|\mathcal{N}(v)|^2] \leq k^2 + k$ for Poisson$(k)$). Thus:
\begin{equation}
    \mathbb{P}(\mathcal{B}) \leq \frac{k^2}{2} \cdot \frac{4k^2}{n} = \frac{2k^4}{n}
\end{equation}

We conclude $\mathbb{P}(\mathcal{B}) = O(k^4/n)$ with explicit constant $C = 8$ (worst-case) or $C = 2$ (expected).

\textbf{Step 3 (Coupling construction):} We construct a coupling $(Y^{\mathrm{joint}}, Y^{\mathrm{prod}})$ as follows:
\begin{enumerate}
    \item Generate the graph $G$ and condition on $Y_v$.
    \item If $\mathcal{B}^c$ holds (no shared neighbors among $\{u_i\}$):
    \begin{itemize}
        \item The labels $\{Y_{u_i}\}$ are conditionally independent given $Y_v$ by the Markov property on trees.
        \item Set $Y^{\mathrm{joint}} = Y^{\mathrm{prod}}$ (perfect coupling).
    \end{itemize}
    \item If $\mathcal{B}$ holds:
    \begin{itemize}
        \item Sample $Y^{\mathrm{joint}}$ from $P_{\mathrm{joint}}$ and $Y^{\mathrm{prod}}$ independently from $P_{\mathrm{prod}}$.
    \end{itemize}
\end{enumerate}

\textbf{Step 4 (TV bound via coupling):} By the coupling characterization of total variation:
\begin{equation}
    d_{\mathrm{TV}}(P_{\mathrm{joint}}, P_{\mathrm{prod}}) \leq \mathbb{P}(Y^{\mathrm{joint}} \neq Y^{\mathrm{prod}}) \leq \mathbb{P}(\mathcal{B}) = O(k^4/n)
\end{equation}

\textbf{Step 5 (Final bound):} Combining the above, under assumption (R1) where $k = O(1)$ is bounded:
\begin{equation}
    d_{\mathrm{TV}}(P_{\mathrm{joint}}, P_{\mathrm{prod}}) \leq \mathbb{P}(\mathcal{B}) \leq \frac{C k^4}{n}
\end{equation}
where $C \leq 8$ (or $C \leq 2$ in expectation). Since $k = O(1)$, this is $O(1/n)$ as $n \to \infty$.
\end{proof}

\begin{remark}[Coupling vs.\ Pinsker]
The coupling proof is preferable because:
\begin{enumerate}
    \item It provides a \textbf{constructive} bound via an explicit bad event.
    \item It avoids the looseness of Pinsker's inequality (which squares the TV distance).
    \item It directly connects to the local tree structure of sparse graphs.
\end{enumerate}
\end{remark}

\begin{corollary}[Error in SNR from Independence Violation]
\label{cor:A2_snr_error}
The violation of conditional independence (A2) introduces an error in $\mathcal{D}(A_v)$ of order:
\begin{equation}
    \left|\mathcal{D}(A_v)_{\mathrm{true}} - \mathcal{D}(A_v)_{\mathrm{indep}}\right| \leq O\left(\frac{k^2}{n}\right) \cdot \mathcal{D}(A_v)_{\mathrm{indep}}
\end{equation}
\end{corollary}

\subsubsection{Multi-Layer Noise Accumulation}

\begin{remark}[Aggregation Normalization Convention]
\label{rem:normalization_convention}
We distinguish two aggregation conventions that lead to different recursions:
\begin{enumerate}
    \item \textbf{Sum aggregation} (used in non-backtracking belief propagation):
    \begin{equation}
        m_{u \to v}^{(t+1)} = X_u + \sum_{w \in \mathcal{N}(u) \setminus \{v\}} m_{w \to u}^{(t)}
    \end{equation}
    Signal grows as $(k-1)^t$, variance grows as $(k-1)^t$, SNR ratio $((k-1)\rho^2)^t$.

    \item \textbf{Mean aggregation} (used in standard GCN):
    \begin{equation}
        h_v^{(t+1)} = \frac{1}{k} \sum_{u \in \mathcal{N}(v)} h_u^{(t)}
    \end{equation}
    Signal decays as $\rho^t$, variance decays as $k^{-t}$, SNR ratio $(\rho^2 k)^t$.
\end{enumerate}

Theorem~\ref{thm:multilayer_with_noise} below uses \textbf{sum aggregation} on trees, which is the natural convention for belief propagation analysis. The SNR ratio $((k-1)\rho^2)^t$ determines whether information amplifies or decays.
\end{remark}

\begin{theorem}[Complete Multi-Layer Recursion with Noise]
\label{thm:multilayer_with_noise}
Consider $t$-layer \textbf{non-backtracking sum aggregation} on the limiting Galton-Watson tree with branching factor $k$. Define:
\begin{itemize}
    \item $\Delta^{(t)} := \|\mathbb{E}[A^{(t)} \mid Y_v = +1] - \mathbb{E}[A^{(t)} \mid Y_v = -1]\|$: Signal magnitude at layer $t$
    \item $V^{(t)} := \mathrm{tr}(\mathrm{Cov}(A^{(t)} \mid Y_v))$: Noise magnitude at layer $t$
\end{itemize}

\textbf{(a) Signal recursion:}
\begin{equation}
    \Delta^{(t+1)} = (k-1) |\rho| \cdot \Delta^{(t)}
\end{equation}

\textbf{(b) Noise recursion:}
\begin{equation}
    V^{(t+1)} = (k-1) V^{(t)} + d_x \sigma_{\varepsilon}^2
\end{equation}
where $\sigma_{\varepsilon}^2$ is per-layer noise variance (from nonlinearity, finite precision, etc.).

\textbf{(c) SNR recursion:}
\begin{equation}
    \mathrm{SNR}^{(t)} = \frac{(\Delta^{(t)})^2}{V^{(t)}} = \frac{((k-1)\rho^2)^t (\Delta^{(0)})^2}{(k-1)^t V^{(0)} + d_x \sigma_{\varepsilon}^2 \frac{(k-1)^t - 1}{k-2}}
\end{equation}

\textbf{(d) Critical layer depth:}
When $\sigma_{\varepsilon}^2 > 0$, there exists a finite optimal depth:
\begin{equation}
    t^* = \arg\max_{t \in \mathbb{N}} \mathrm{SNR}^{(t)}
\end{equation}
Beyond $t^*$, noise accumulation dominates and SNR decreases (oversmoothing).
\end{theorem}

\begin{proof}
\textbf{Part (a):} By the non-backtracking structure, each message aggregates over $k-1$ new branches. The expected signal scales by $\rho$ per edge, giving $(k-1)\rho$ per layer.

\textbf{Part (b):} The noise from $k-1$ children adds up. If each child contributes variance $V^{(t)}$ and there is additional per-layer noise $\sigma_{\varepsilon}^2$:
\begin{equation}
    V^{(t+1)} = (k-1) V^{(t)} + d_x \sigma_{\varepsilon}^2
\end{equation}

\textbf{Part (c):} Solving the linear recursion for $V^{(t)}$:
\begin{equation}
    V^{(t)} = (k-1)^t V^{(0)} + d_x \sigma_{\varepsilon}^2 \sum_{s=0}^{t-1} (k-1)^s = (k-1)^t V^{(0)} + d_x \sigma_{\varepsilon}^2 \frac{(k-1)^t - 1}{k-2}
\end{equation}

\textbf{Part (d):} Taking derivative of $\log \mathrm{SNR}^{(t)}$ with respect to $t$ and setting to zero gives the optimal depth. For large $t$, the noise term dominates and SNR decays.
\end{proof}

\subsubsection{Unified Error Bound for Main Theorem}

\begin{theorem}[Complete Discriminability Formula with Error]
\label{thm:complete_with_error}
Under the sparse CSBM (Definition~\ref{def:sparse_csbm}) with:
\begin{itemize}
    \item $n$ nodes, average degree $k = O(1)$
    \item Feature SNR $\mathcal{D}(X) = 4\mu^\top \Sigma^{-1} \mu$
    \item $t$-layer non-backtracking aggregation with $t = o(\log n)$
\end{itemize}

The discriminability satisfies:
\begin{equation}
    \mathcal{D}^{(t)} = ((k-1)\rho^2)^t \cdot \mathcal{D}(X) \cdot (1 + \varepsilon_n(t))
\end{equation}
where the error term is:
\begin{equation}
    |\varepsilon_n(t)| \leq \underbrace{\frac{C_1 k^2}{n}}_{\text{Independence violation}} + \underbrace{\frac{C_2 k^{2t}}{n}}_{\text{Cycle contamination}} + \underbrace{\frac{C_3 \|\mu\|^2 (1-\rho^2)}{\sigma^2}}_{\text{Label mixing noise}}
\end{equation}
for universal constants $C_1, C_2, C_3 > 0$.
\end{theorem}

\begin{remark}[Validity Regime and Error Decomposition]
\label{rem:validity_regime_error}
The error $\varepsilon_n(t)$ consists of three components with different scaling behaviors:
\begin{enumerate}
    \item \textbf{Graph-dependent terms} ($O(1/n)$): The first two terms $\frac{C_1 k^2}{n} + \frac{C_2 k^{2t}}{n}$ vanish as $n \to \infty$ for fixed $k$ and $t = o(\log_k n)$.

    \item \textbf{Model-dependent term} ($O(1)$): The label mixing term $\frac{C_3 \|\mu\|^2 (1-\rho^2)}{\sigma^2}$ is \textbf{bounded but non-vanishing} under Assumption R4 ($\|\mu\|^2/\sigma^2 = O(1)$). This term captures the intrinsic approximation error from treating neighbor labels as independent.
\end{enumerate}

\textbf{Interpretation:} The theorem provides a multiplicative approximation $\mathcal{D}^{(t)} \approx ((k-1)\rho^2)^t \mathcal{D}(X)$ with:
\begin{itemize}
    \item \textbf{Asymptotic accuracy}: Error from graph structure vanishes as $n \to \infty$.
    \item \textbf{Bounded bias}: A constant multiplicative factor $(1 \pm O(\|\mu\|^2/\sigma^2))$ remains even in the limit.
\end{itemize}

For high feature SNR ($\|\mu\|^2/\sigma^2 \gg 1$), the label mixing term dominates, and the approximation degrades. However, this regime is less interesting practically: when features alone are highly informative, graph structure provides minimal additional benefit regardless of $\kappa$.
\end{remark}

\begin{lemma}[Multi-Layer Error Composition]
\label{lem:error_composition}
Under the sparse CSBM with $t$-layer non-backtracking aggregation, the total error $\varepsilon_n(t)$ decomposes as:
\begin{equation}
    |\varepsilon_n(t)| \leq \varepsilon_{\mathrm{indep}}(t) + \varepsilon_{\mathrm{cycle}}(t) + \varepsilon_{\mathrm{mix}}
\end{equation}
where each component satisfies:
\begin{enumerate}
    \item \textbf{Independence violation}: $\varepsilon_{\mathrm{indep}}(t) = O(t \cdot k^4/n)$ (additive across layers)
    \item \textbf{Cycle contamination}: $\varepsilon_{\mathrm{cycle}}(t) = O((k-1)^{2t}/n)$ (exponential in $t$)
    \item \textbf{Label mixing}: $\varepsilon_{\mathrm{mix}} = O(\|\mu\|^2(1-\rho^2)/\sigma^2)$ (layer-independent)
\end{enumerate}
\end{lemma}

\begin{proof}
\textbf{Part 1 (Independence violation):}
By Lemma~\ref{lem:A2_quantitative}, at each layer $\ell$, the conditional independence violation contributes $d_{\mathrm{TV}} = O(k^4/n)$ to the joint distribution error.

On the limiting Galton-Watson tree, non-backtracking aggregation ensures that the neighbors at layer $\ell$ are disjoint from those at layer $\ell' \neq \ell$ (except for the single backtrack edge). Thus, the TV errors at different layers are \textit{independent} and combine additively:
\begin{equation}
    d_{\mathrm{TV}}^{(1:t)} \leq \sum_{\ell=1}^t d_{\mathrm{TV}}^{(\ell)} = O(t \cdot k^4/n)
\end{equation}

\textbf{Part 2 (Cycle contamination):}
The tree approximation fails when the $t$-hop neighborhood $B_t(v)$ contains a cycle. The expected number of cycles of length $\geq 3$ in $B_t(v)$ is:
\begin{equation}
    \mathbb{E}[\#\text{cycles}] \leq \sum_{s=3}^{2t} \frac{((k-1)^t)^2}{n} \cdot O(1) = O\left(\frac{(k-1)^{2t}}{n}\right)
\end{equation}
By Markov's inequality, $\mathbb{P}(\text{cycle exists}) = O((k-1)^{2t}/n)$.

\textbf{Part 3 (Label mixing):}
The label mixing noise arises from the variance term $\mu\mu^\top \mathrm{Var}(\sum_u w_{vu} Y_u \mid Y_v)$ in the covariance decomposition. Under condition (A3), this contributes a relative error:
\begin{equation}
    \frac{\|\mu\|^2 (1-\rho^2)}{\sigma^2 / k_{\mathrm{eff}}} = O\left(\frac{\|\mu\|^2 (1-\rho^2)}{\sigma^2}\right) \cdot k_{\mathrm{eff}}
\end{equation}
For bounded $k_{\mathrm{eff}} = O(1)$, this is $O(\|\mu\|^2(1-\rho^2)/\sigma^2)$ and does not depend on $t$.
\end{proof}

\begin{remark}[Error Metric Unification]
\label{rem:error_metric}
The three error terms in Lemma~\ref{lem:error_composition} measure different quantities but combine to bound the \textbf{relative error in discriminability} $\mathcal{D}$:
\begin{itemize}
    \item $\varepsilon_{\mathrm{indep}}$ bounds the TV distance between true and idealized neighbor label distributions. Under our Gaussian model, TV error $\delta$ translates to relative error in $\mathcal{D}$ as follows: by Pinsker's inequality $\mathrm{TV}^2 \leq \frac{1}{2} D_{\mathrm{KL}}$, and for Gaussians differing in mean by $\Delta\mu$, $D_{\mathrm{KL}} = O(\|\Delta\mu\|^2/\sigma^2)$. Since $\mathcal{D} \propto \|\Delta\|^2/\sigma^2$, perturbations in mean of order $\delta\|\mu\|$ yield $O(\delta)$ relative error in $\mathcal{D}$.
    \item $\varepsilon_{\mathrm{cycle}}$ bounds the probability that the local neighborhood deviates from a tree. On the complement event (tree-like), the SNR recursion holds exactly; by the law of total expectation, conditioning contributes $O(\mathbb{P}(\text{cycle}))$ additive error to any bounded statistic.
    \item $\varepsilon_{\mathrm{mix}}$ is directly a relative error in $\mathcal{D}$ from the variance decomposition (Theorem~\ref{thm:agg_discriminability_rigorous}(b)).
\end{itemize}
Thus, the total relative error $|\mathcal{D}_{\mathrm{true}} - \mathcal{D}_{\mathrm{ideal}}|/\mathcal{D}_{\mathrm{ideal}} \leq \varepsilon_n(t)$ is bounded by their sum.

\textbf{Regularity conditions for TV$\to$discriminability bridge:} The above arguments require:
\begin{enumerate}
    \item[(i)] \textbf{Bounded signal-to-noise}: $\|\mu\|^2/\sigma^2 = O(1)$ (Assumption R4), ensuring Gaussian tails are well-behaved.
    \item[(ii)] \textbf{Non-degenerate covariance}: $\|\Sigma^{-1}\|_{\mathrm{op}} = O(1)$, so that small perturbations in distribution translate to controlled perturbations in $\mathcal{D}$.
    \item[(iii)] \textbf{Explicit Lipschitz bound}: The discriminability $\mathcal{D}(\mu, \Sigma) = 4\mu^\top\Sigma^{-1}\mu$ satisfies, for perturbations $(\Delta\mu, \Delta\Sigma)$:
    \begin{equation}
        |\mathcal{D}(\mu + \Delta\mu, \Sigma) - \mathcal{D}(\mu, \Sigma)| \leq 8\|\Sigma^{-1}\|_{\mathrm{op}} \|\mu\| \|\Delta\mu\|
    \end{equation}
    (by direct differentiation: $\nabla_\mu \mathcal{D} = 8\Sigma^{-1}\mu$). Under conditions (i)-(ii), the Lipschitz constant $L_\mu = 8\|\Sigma^{-1}\|_{\mathrm{op}} \|\mu\| = O(1)$ is bounded uniformly in $n$.
\end{enumerate}

\textbf{Source of perturbations and TV$\to$mean bridge:} In our framework, mean perturbations arise from the TV distance $\delta$ between true and idealized neighbor label distributions. The key structural property enabling the TV$\to$mean bridge is:

\begin{lemma}[TV to Mean Perturbation under Linear Model]
\label{lem:tv_to_mean}
Let $Y^{\mathrm{true}} = (Y_{u_1}, \ldots, Y_{u_k})$ and $Y^{\mathrm{ideal}}$ be the true and idealized neighbor label vectors with $d_{\mathrm{TV}}(Y^{\mathrm{true}}, Y^{\mathrm{ideal}}) \leq \delta$, where we use the standard definition $d_{\mathrm{TV}}(P, Q) := \sup_{A} |P(A) - Q(A)| = \frac{1}{2}\|P - Q\|_1$. Under the linear feature model $X_u = Y_u \mu + \varepsilon_u$ with $Y_u \in \{-1, +1\}$ and \textbf{deterministic} normalized aggregation weights $w_{vu} \geq 0$, $\sum_u w_{vu} = 1$, the aggregated mean satisfies:
\begin{equation}
    \|\mathbb{E}[A_v]_{\mathrm{true}} - \mathbb{E}[A_v]_{\mathrm{ideal}}\| \leq 2\delta \|\mu\|
\end{equation}

\textbf{Note on weight conditioning:} The weights $\{w_{vu}\}$ are treated as \textit{fixed} (deterministic) when applying this lemma. In practice, weights may depend on the graph structure $G$ (e.g., $w_{vu} = 1/d_v$ for mean aggregation). In such cases, the lemma is applied \textit{conditional on} the realized graph $G$, and the bound holds almost surely over $G$. Since the TV bound in Lemma~\ref{lem:A2_quantitative} is already conditional on the neighborhood structure $\mathcal{N}(v)$, the conditioning is consistent throughout our analysis.
\end{lemma}

\begin{proof}
The aggregated feature is $A_v = \sum_u w_{vu} X_u = \sum_u w_{vu} (Y_u \mu + \varepsilon_u)$. Taking expectations and using linearity:
\begin{equation}
    \mathbb{E}[A_v \mid Y_v] = \mu \cdot \underbrace{\sum_u w_{vu} \mathbb{E}[Y_u \mid Y_v]}_{=: f(Y) \text{, a scalar}}
\end{equation}
Thus the mean perturbation is $\|\mathbb{E}[A_v]_{\mathrm{true}} - \mathbb{E}[A_v]_{\mathrm{ideal}}\| = \|\mu\| \cdot |f_{\mathrm{true}} - f_{\mathrm{ideal}}|$ where $f = \sum_u w_{vu} Y_u$.

For the scalar function $f: \{-1,+1\}^k \to \mathbb{R}$, we have the bound $|f| \leq \sum_u w_{vu} |Y_u| = \sum_u w_{vu} = 1$ (since $Y_u \in \{-1,+1\}$ and weights are non-negative summing to 1).

By the standard TV-expectation inequality (Tsybakov~\cite{tsybakov2009introduction}, Lemma~2.1): for any \textbf{measurable} function $f: \mathcal{X} \to \mathbb{R}$ with $\|f\|_\infty \leq M$,
\begin{equation}
    |\mathbb{E}_{P}[f] - \mathbb{E}_{Q}[f]| \leq 2M \cdot d_{\mathrm{TV}}(P, Q)
\end{equation}
The function $f(Y) = \sum_u w_{vu} Y_u$ is measurable (it is a finite linear combination of coordinate projections on the discrete space $\{-1,+1\}^k$).
Applying this with $M = 1$:
\begin{equation}
    |f_{\mathrm{true}} - f_{\mathrm{ideal}}| = \left|\sum_u w_{vu} \mathbb{E}[Y_u]_{\mathrm{true}} - \sum_u w_{vu} \mathbb{E}[Y_u]_{\mathrm{ideal}}\right| \leq 2\delta
\end{equation}
Multiplying by $\|\mu\|$ gives the result.
\end{proof}

Combining with the Lipschitz bound (condition (iii)), if $d_{\mathrm{TV}}(P_{\mathrm{joint}}, P_{\mathrm{prod}}) \leq \delta$, then:
\begin{equation}
    |\mathcal{D}_{\mathrm{true}} - \mathcal{D}_{\mathrm{ideal}}| \leq L_\mu \cdot 2\delta \|\mu\| = 16\|\Sigma^{-1}\|_{\mathrm{op}} \|\mu\|^2 \delta = O(\delta) \cdot \mathcal{D}_{\mathrm{ideal}}
\end{equation}
where we used the spectral bound $\mathcal{D}_{\mathrm{ideal}} = 4\mu^\top \Sigma^{-1} \mu \geq 4\|\mu\|^2 / \|\Sigma\|_{\mathrm{op}}$ (since $\Sigma^{-1} \succeq \|\Sigma\|_{\mathrm{op}}^{-1} I$). The relative error is thus bounded by $4\kappa(\Sigma) \delta$ where $\kappa(\Sigma) = \|\Sigma\|_{\mathrm{op}} \|\Sigma^{-1}\|_{\mathrm{op}}$ is the condition number, which is $O(1)$ under our assumptions.

\textbf{Lipschitz bound locality:} The Lipschitz constant $L_\mu = 8\|\Sigma^{-1}\|_{\mathrm{op}} \|\mu\|$ depends on the reference point $\mu$. This is appropriate since $\mathcal{D}(\mu, \Sigma) = 4\mu^\top \Sigma^{-1} \mu$ is quadratic in $\mu$. Under Assumption (R4) where $\|\mu\|^2/\sigma^2 = O(1)$, the parameter $\mu$ is fixed (independent of $n$), so $L_\mu = O(1)$ uniformly.

\textbf{Covariance invariance:} The covariance $\Sigma$ refers to the \textit{conditional} variance $\mathrm{Var}(X_u \mid Y_u) = \mathrm{Var}(\varepsilon_u) = \Sigma$, which remains unchanged ($\Delta\Sigma = 0$) because:
\begin{enumerate}
    \item The feature noise $\varepsilon_u \sim \mathcal{N}(0, \Sigma)$ is i.i.d.\ and independent of the graph structure.
    \item Our coupling construction (Lemma~\ref{lem:A2_quantitative}) only modifies the \textit{joint distribution of neighbor labels}, not the noise model.
    \item Cycle contamination affects correlations between neighbor labels, not the conditional variance of $X_u \mid Y_u$.
\end{enumerate}

These conditions are satisfied under our sparse CSBM (Definition~\ref{def:sparse_csbm}) with fixed parameters $(\mu, \Sigma, a, b)$ independent of $n$.
\end{remark}

\begin{remark}[Layer-Independence of $\varepsilon_{\mathrm{mix}}$]
\label{rem:mix_layer_indep}
The label mixing term $\varepsilon_{\mathrm{mix}} = O(\|\mu\|^2(1-\rho^2)/\sigma^2)$ does \textbf{not} accumulate across layers because:
\begin{enumerate}
    \item \textbf{Modeling assumption}: In our framework, stochastic noise is introduced \textit{only at feature generation} ($X_v = Y_v\mu + \varepsilon_v$). The aggregation operations are \textit{deterministic linear transformations} given the graph and neighbor features.
    \item \textbf{No per-layer noise injection}: Unlike some practical GNN implementations with dropout or stochastic regularization, our theoretical model assumes clean message passing. Any ``noise'' in aggregated representations comes solely from the initial feature noise $\{\varepsilon_v\}$.
    \item \textbf{Relative error interpretation}: The mixing term captures how much the neighbor label uncertainty (controlled by $\rho$) contaminates the signal direction $\mu$. This contamination is determined by the \textit{one-hop} edge correlation structure, not the depth of propagation.
\end{enumerate}
If per-layer noise were present, one would need $\varepsilon_{\mathrm{mix}}(t) = O(t \cdot \|\mu\|^2(1-\rho^2)/\sigma^2)$ or a geometric series bound.
\end{remark}

\begin{remark}[Error Accumulation Across Layers]
\label{rem:error_accumulation}
The key insight from Lemma~\ref{lem:error_composition} is that different error sources have different accumulation behaviors:
\begin{itemize}
    \item \textbf{Independence violation}: Additive in $t$, giving $O(t \cdot k^4/n)$
    \item \textbf{Cycle contamination}: Exponential in $t$, giving $O((k-1)^{2t}/n)$ --- this is the bottleneck
    \item \textbf{Label mixing}: Layer-independent, bounded by feature SNR
\end{itemize}

\textbf{Unified condition for validity:} All three error sources are $o(1)$ when:
\begin{equation}
    t \leq \left(\frac{1}{2} - \epsilon\right) \log_{k-1} n \quad \text{(ensures } (k-1)^{2t}/n \to 0)
\end{equation}
which is precisely condition (R2) in Assumption~\ref{ass:validity_regime}.
\end{remark}

\begin{remark}[Single-Node vs.\ Whole-Graph Applicability]
\label{rem:single_vs_whole}
Our theoretical results are stated for a \textbf{fixed root node $v$} and characterize the discriminability $\mathcal{D}(A_v)$ of its aggregated representation. We clarify the scope:

\textbf{(a) Local analysis (our primary setting):} The bounds in Theorem~\ref{thm:complete_with_error} hold for any \textit{fixed} node $v$ with probability $1 - o(1)$ over the random graph. This is the natural setting for understanding how aggregation affects a typical node's classification.

\textbf{(b) Whole-graph uniformity:} To extend the results to \textit{all} nodes simultaneously, one applies a union bound.

\textbf{Precise single-node failure probability:} From Lemma~\ref{lem:error_composition}, the per-node failure probability (event that tree approximation fails) is:
\begin{equation}
    p_{\mathrm{fail}}(v, t) := \mathbb{P}(B_t(v) \text{ contains a cycle}) = \frac{C_{\mathrm{cycle}} (k-1)^{2t}}{n}
\end{equation}
where $C_{\mathrm{cycle}} > 0$ is a \textbf{universal constant for all $v \in V$}, depending only on $(a, b)$ (not on $n$, $t$, or the specific node $v$). This uniformity follows from:
\begin{enumerate}
    \item The CSBM is vertex-exchangeable: the law of $B_t(v)$ is identical for all $v$ before conditioning on labels.
    \item The expected number of cycles formula $\mathbb{E}[\#\text{cycles in } B_t(v)] = \sum_{s=3}^{2t} O((k-1)^{2t}/n)$ involves only global parameters $(a, b, n)$.
    \item The constant $C_{\mathrm{cycle}} \leq \sum_{s=3}^{\infty} c_s$ where $c_s$ depends only on the cycle counting combinatorics and $(a+b)^s/n^{s-1}$.
\end{enumerate}

\textbf{Union bound for all nodes:}
\begin{equation}
    \mathbb{P}\left(\exists v \in V: |\mathcal{D}(A_v) - \mathcal{D}_{\mathrm{ideal}}| > \delta \cdot \mathcal{D}_{\mathrm{ideal}}\right) \leq n \cdot p_{\mathrm{fail}}(v, t) = C_{\mathrm{cycle}} (k-1)^{2t}
\end{equation}
For this to vanish as $n \to \infty$, we need $(k-1)^{2t} \to 0$, which requires $t = O(1)$ (bounded depth) for \textbf{exact uniformity over all vertices}.

\textbf{Relaxation to ``most vertices'':} If we only require the fraction of ``bad'' vertices to vanish, i.e., $\mathbb{E}[|\{v: \text{bad}\}|]/n \to 0$, then by Markov's inequality:
\begin{equation}
    \frac{\mathbb{E}[|\{v: \text{bad}\}|]}{n} = p_{\mathrm{fail}}(v, t) = \frac{C_{\mathrm{cycle}} (k-1)^{2t}}{n} \to 0
\end{equation}
holds for any $t = o(\log_{k-1} \sqrt{n})$, which is much weaker than the per-node constraint $t = o(\log_{k-1} n)$ but still allows $t$ to grow with $n$.

\textbf{Slowly growing $t$ with high probability:} For any $\omega(n) \to \infty$ arbitrarily slowly, taking:
\begin{equation}
    t \leq \frac{\log\log n - \log\omega(n)}{2\log(k-1)}
\end{equation}
ensures $(k-1)^{2t} \leq (\log n)/\omega(n) \to 0$, hence uniform validity with probability $1 - o(1)$.

\textbf{(c) Practical implication:} For typical GNN depths ($t \leq 3$) and sparse graphs ($k = O(1)$), both per-node and whole-graph analyses yield the same asymptotic behavior. The distinction matters only for deep propagation ($t = \Omega(\log n)$), which is rare in practice due to oversmoothing.

\textbf{(d) Empirical measurements:} Our experimental results (Section~\ref{sec:experiments}) report \textit{average} accuracy over all nodes, which corresponds to the expected discriminability $\mathbb{E}_v[\mathcal{D}(A_v)]$. By linearity of expectation and the sparse limit theorem, this equals the per-node analysis up to $o(1)$ error.
\end{remark}

\subsubsection{Unified Probability Space for Error Control}

\begin{remark}[Unified Probability Space]
\label{rem:unified_probability_space}
We clarify the probability space underlying all error bounds in this paper, addressing potential confusion about different sources of randomness.

\textbf{(a) Fundamental probability space:}
Let $(\Omega, \mathcal{F}, \mathbb{P})$ be the product probability space where:
\begin{itemize}
    \item \textbf{Labels}: $\sigma = (Y_1, \ldots, Y_n) \in \{-1, +1\}^n$ with $Y_v \stackrel{\text{iid}}{\sim} \mathrm{Rad}(1/2)$
    \item \textbf{Graph}: $G \in \{0,1\}^{n \times n}$ with edges sampled conditionally on $\sigma$ per CSBM
    \item \textbf{Features}: $X = (X_1, \ldots, X_n) \in \mathbb{R}^{n \times d_x}$ with $X_v = Y_v \mu + \varepsilon_v$, $\varepsilon_v \stackrel{\text{iid}}{\sim} \mathcal{N}(0, \Sigma)$
\end{itemize}
All three components are generated on a \textit{single} probability space with the factorization:
\begin{equation}
    \mathbb{P}(\sigma, G, X) = \mathbb{P}(\sigma) \cdot \mathbb{P}(G \mid \sigma) \cdot \mathbb{P}(X \mid \sigma)
\end{equation}
Note that $X$ and $G$ are conditionally independent given $\sigma$.

\textbf{(b) Error bounds hold jointly:}
All error bounds in this paper (Lemmas~\ref{lem:R_n_bound}, \ref{lem:A2_quantitative}, Theorem~\ref{thm:complete_with_error}) are stated as:
\begin{equation}
    \mathbb{P}_{\sigma, G, X}\left(\text{``bad event''}\right) \leq \delta(n)
\end{equation}
where $\delta(n) \to 0$ as $n \to \infty$. By a union bound, all error events can be controlled simultaneously:
\begin{equation}
    \mathbb{P}\left(\bigcup_{i=1}^{m} E_i\right) \leq \sum_{i=1}^{m} \mathbb{P}(E_i) = o(1)
\end{equation}
for any finite number $m$ of error events.

\textbf{(c) Conditioning hierarchy:}
Our analysis uses the following conditioning structure:
\begin{enumerate}
    \item \textbf{Condition on $\sigma$}: The discriminability $\mathcal{D}(A_v)$ is defined for a \textit{fixed} label configuration $\sigma$. The randomness in $\mathcal{D}(A_v)$ comes from $G$ and $X$.
    \item \textbf{Condition on $\sigma$ and $G$}: The TV bound (Lemma~\ref{lem:A2_quantitative}) holds for fixed $(\sigma, G)$. The weights $\{w_{vu}\}$ are then deterministic.
    \item \textbf{Condition on $\sigma$, $G$, and neighbors}: The mean/variance computations (Theorem~\ref{thm:agg_discriminability_rigorous}) use the conditional distribution $X_u \mid Y_u$.
\end{enumerate}

\textbf{(d) Asymptotic notation clarification:}
When we write $\mathcal{D}^{(t)} = ((k-1)\rho^2)^t \mathcal{D}^{(0)} \cdot (1 + \varepsilon_n(t))$ with $|\varepsilon_n(t)| = o(1)$, this means:
\begin{equation}
    \forall \delta > 0, \exists N(\delta) \text{ s.t. } n \geq N(\delta) \Rightarrow \mathbb{P}(|\varepsilon_n(t)| > \delta) \leq \delta
\end{equation}
That is, convergence is in probability over the joint space $(\sigma, G, X)$.

\textbf{(e) Non-asymptotic alternative:}
For finite $n$, all bounds can be made non-asymptotic by tracking constants explicitly. For instance, Theorem~\ref{thm:complete_with_error} holds with:
\begin{equation}
    \mathbb{P}\left(|\varepsilon_n(t)| > C\left(\frac{k^4}{n} + \frac{k^{2t}}{n} + \frac{\|\mu\|^2(1-\rho^2)}{\sigma^2}\right)\right) \leq 2e^{-n/C'}
\end{equation}
for explicit universal constants $C, C' > 0$ depending only on model parameters $(a, b, d_x)$.
\end{remark}

\begin{remark}[Independence of Error Sources]
\label{rem:error_independence}
The three error sources in Lemma~\ref{lem:error_composition} are \textbf{not independent} in general, but they are \textbf{separately controllable}:
\begin{enumerate}
    \item $\varepsilon_{\mathrm{indep}}$ depends on graph structure (presence of shared neighbors)
    \item $\varepsilon_{\mathrm{cycle}}$ depends on graph structure (presence of cycles)
    \item $\varepsilon_{\mathrm{mix}}$ depends on model parameters $(\mu, \sigma, \rho)$, not graph realization
\end{enumerate}

The key insight is that $\varepsilon_{\mathrm{mix}}$ is \textbf{deterministic} given model parameters, while $\varepsilon_{\mathrm{indep}}$ and $\varepsilon_{\mathrm{cycle}}$ are random but \textbf{positively correlated} (both are worse when the graph has denser local structure). Thus:
\begin{equation}
    \mathbb{P}(|\varepsilon_n(t)| > \eta) \leq \mathbb{P}(\varepsilon_{\mathrm{indep}} + \varepsilon_{\mathrm{cycle}} > \eta - \varepsilon_{\mathrm{mix}})
\end{equation}
The right-hand side can be bounded using concentration for graph functionals (e.g., Janson's inequality for cycle counts), giving sub-Gaussian tails.
\end{remark}

\subsubsection{Bayes Error Consistency}

\begin{proposition}[Bayes Error in Terms of Discriminability]
\label{prop:bayes_error_consistent}
Under the Gaussian model, the Bayes error satisfies:
\begin{equation}
    P_e = \Phi\left(-\frac{1}{2}\sqrt{\mathcal{D}}\right) = \Phi\left(-\frac{1}{2}\sqrt{\Delta^\top \Sigma^{-1} \Delta}\right)
\end{equation}

For sub-Gaussian noise with parameter $\sigma$, the error of the linear classifier (projection onto signal direction) is bounded:
\begin{equation}
    P_e^{\mathrm{linear}} \leq \exp\left(-\frac{\mathcal{D}}{8}\right)
\end{equation}
when $\mathcal{D} \geq 2$.

The connection to our SNR notation:
\begin{equation}
    \mathcal{D}(A) = \kappa \cdot \mathcal{D}(X) \quad \Rightarrow \quad P_e(A) \approx \Phi\left(-\frac{1}{2}\sqrt{\kappa \cdot \mathcal{D}(X)}\right)
\end{equation}
\end{proposition}

\begin{proof}
The Gaussian result is standard (optimal linear discriminant). For sub-Gaussian, we use Hoeffding's inequality on the projected statistic:
\begin{equation}
    T = \mu^\top \Sigma^{-1} X \sim \text{sub-Gaussian with variance proxy } \mu^\top \Sigma^{-1} \mu = \mathcal{D}/4
\end{equation}
Under $Y = +1$, $\mathbb{E}[T] = \mathcal{D}/2$. Thus:
\begin{equation}
    P(T < 0 \mid Y = +1) \leq \exp\left(-\frac{(\mathcal{D}/2)^2}{2 \cdot \mathcal{D}/4}\right) = \exp\left(-\frac{\mathcal{D}}{8}\right)
\end{equation}
\end{proof}

\subsubsection{Definition of $\rho$ from SBM Parameters}

\begin{definition}[Correlation Coefficient $\rho$]
\label{def:rho_explicit}
In the sparse CSBM with within-class edge probability $a/n$ and between-class edge probability $b/n$:
\begin{equation}
    \rho := \mathbb{E}[Y_u \mid (u,v) \in E, Y_v = +1] = \frac{a - b}{a + b}
\end{equation}

This can be rewritten in terms of homophily $h = a/(a+b)$:
\begin{equation}
    \rho = 2h - 1
\end{equation}

\textbf{Estimator:} Given observed edges and labels:
\begin{equation}
    \hat{\rho} = \frac{E_{\mathrm{in}} - E_{\mathrm{out}}}{E_{\mathrm{in}} + E_{\mathrm{out}}}
\end{equation}
where $E_{\mathrm{in}} = |\{(u,v) \in E : Y_u = Y_v\}|$ counts within-class edges.

\textbf{Under GCN normalization:} The effective correlation becomes:
\begin{equation}
    \rho_w = \sum_{j \in \mathcal{N}(i)} w_{ij} \cdot \mathbb{E}[Y_j \mid Y_i = +1]
\end{equation}
For symmetric normalization $w_{ij} = 1/\sqrt{d_i d_j}$ and approximately regular graphs, $\rho_w \approx \rho$.
\end{definition}

\subsection{Summary: Complete Rigorous Framework}
\label{subsec:complete_summary_rigorous}

\begin{table}[t]
\centering
\caption{Complete SNR Framework with Quantitative Errors}
\label{tab:complete_rigorous_summary}
\begin{tabular}{@{}lccc@{}}
\toprule
\textbf{Result} & \textbf{Formula} & \textbf{Error Bound} & \textbf{Conditions} \\
\midrule
\multicolumn{4}{l}{\textit{Core Results}} \\
Discriminability & $\mathcal{D} = \Delta^\top \Sigma^{-1} \Delta$ & Exact & Definition \\
Bayes error & $P_e = \Phi(-\frac{1}{2}\sqrt{\mathcal{D}})$ & Exact & Gaussian \\
One-hop $\kappa$ & $\rho^2 k_{\mathrm{eff}}$ & $O(k^2/n + \|\mu\|^2/\sigma^2)$ & Sparse limit \\
\midrule
\multicolumn{4}{l}{\textit{Multi-Layer}} \\
$t$-layer SNR & $((k-1)\rho^2)^t \mathcal{D}^{(0)}$ & $O(k^{2t}/n)$ & $t = o(\log n)$ \\
With noise & Full recursion (Thm~\ref{thm:multilayer_with_noise}) & --- & Tree regime \\
\midrule
\multicolumn{4}{l}{\textit{Error Sources}} \\
$R_n$ bound & $\|R_n\|_{\mathrm{op}} = O(k^2/n)$ & --- & Spectral norm \\
A2 violation & $d_{\mathrm{TV}} = O(k^4/n)$ & Coupling & Sparse limit \\
Label mixing & $O(\|\mu\|^2(1-\rho^2)/\sigma^2)$ & --- & Relative scale \\
\bottomrule
\end{tabular}
\end{table}

With these quantitative bounds, the theoretical framework is now complete and defensible at the TKDE level. The key improvements are:
\begin{enumerate}
    \item Explicit error terms $\varepsilon_n(t)$ with quantified dependence on $n, k, t, \|\mu\|, \sigma$
    \item Total variation bound via coupling for conditional independence (Lemma~\ref{lem:A2_quantitative})
    \item Complete multi-layer recursion including noise accumulation (Theorem~\ref{thm:multilayer_with_noise})
    \item Consistent definition of $\rho$ from SBM parameters with estimator
    \item Unified Bayes error formula compatible with $\mathcal{D}$ definition
    \item Explicit error accumulation across layers (Remark~\ref{rem:error_accumulation})
    \item Unified notation for $k$, $k-1$, and $k_{\mathrm{eff}}$ (Remark~\ref{rem:k_notation_unified})
\end{enumerate}


